\documentclass[10pt]{article}

\usepackage[margin=0.82in]{geometry}
\usepackage[T1]{fontenc}
\usepackage[utf8]{inputenc}
\usepackage{lmodern}
\usepackage{microtype}
\usepackage{amsmath,amssymb,mathtools}
\usepackage{booktabs}
\usepackage{tabularx}
\usepackage{array}
\usepackage{enumitem}
\usepackage{xurl}
\usepackage[colorlinks=true,allcolors=blue]{hyperref}

\setlength{\parskip}{0.35em}
\setlength{\parindent}{0pt}
\setlist{nosep,leftmargin=*}
\renewcommand{\arraystretch}{1.12}
\newcolumntype{Y}{>{\raggedright\arraybackslash}X}
\newcommand{\method}{\textsc{CCD-GS}}
\newcommand{\pending}{\textit{Pending measurement}}

\title{\textbf{Do Temporary Pose Probes Predict Useful Gaussian Splits?}\\
       \large A Split-Utility Study for 3D Gaussian Splatting}
\author{\textit{Author information to be added}}
\date{Working draft --- 10 October 2026}

\begin{document}
\maketitle

\begin{center}
\fbox{\parbox{0.94\linewidth}{
\textbf{Draft status.} This manuscript records a completed Stage 0 research
specification and data-preparation audit. No \method{} training, pose-probe,
split-utility, end-to-end allocation, or external-baseline experiment has been
run. All method benefits and empirical conclusions are therefore unestablished.
The Results and Discussion sections record this evidence boundary rather than
inventing outcomes.
}}
\end{center}

\begin{abstract}
Adaptive density control in 3D Gaussian Splatting (3DGS) uses image-space
optimization signals to decide where the representation should grow. Such
signals can also change when camera poses are refined, so a large growth
gradient need not uniquely indicate that an additional Gaussian will improve
the reconstruction. We specify a controlled study of whether a bounded,
temporary intervention on camera extrinsics provides information about the
subsequent utility of splitting an eligible Gaussian. The intervention fits a
copied camera while holding the current scene values fixed, measures the same
screen-space growth signal before and after the fit, and discards the temporary
correction. Its predictive value is to be tested with matched split/no-split
training branches and compared with original-gradient, corrected-gradient,
retained-correction, and equal-computation pose-refinement controls. The study
prioritizes reconstruction quality versus Gaussian count and separately
accounts for training time, probe overhead, memory, and camera error. We have
prepared a known-camera synthetic fixture and an official low-resolution
Mip-NeRF 360 room scene with explicit training, development, and reserved-test
view manifests. These preparation checks do not constitute model experiments.
No conclusion about split prediction, compactness, speed, or novelty is claimed
in this draft.
\end{abstract}

\textbf{Keywords:} 3D Gaussian Splatting, density control, camera pose
refinement, split utility, representation allocation

\section{Introduction}

3D Gaussian Splatting represents a scene with explicit anisotropic Gaussian
primitives and renders them through a differentiable rasterization pipeline
\cite{kerbl2023}. Its adaptive density-control procedure can clone or split
primitives as training proceeds. These operations determine where model
capacity is added and affect the final primitive population. A split can
improve a region that the current representation cannot fit, but the
optimization signal used to nominate a candidate is not itself a measurement
of the benefit that the split will produce.

One source of ambiguity is camera pose. A change in extrinsics alters
projected Gaussian means, visibility, and image loss. Consequently, the
screen-space gradient of an eligible Gaussian may respond to pose refinement
even though its scene parameters have not changed. This response is
potentially useful when deciding where to allocate more primitives, but it has
several alternative explanations. A camera update can compensate for incorrect
geometry; it can change occlusion; and a gradient that persists after pose
refinement can still be resolved by optimizing existing Gaussians rather than
adding capacity. A gradient decrease alone therefore does not show that a
split was unnecessary.

This paper asks a narrower question:

\begin{quote}
\textbf{Does a bounded, temporary pose intervention identify useful Gaussian
splits better than original or corrected gradients, and does that information
remain useful after accounting for equivalent permanent camera refinement?}
\end{quote}

We use \method{} as a project name for the proposed investigation. At a
training checkpoint, the diagnostic copies one training camera, holds all
Gaussian values fixed, fits a bounded extrinsic correction using that training
image, measures a candidate's normalized screen-space growth signal again,
and discards the correction. The measurements are treated as diagnostic
features, not as proof of calibration error or as split-utility labels. The
central mechanism test instead restores identical training states into
matched branches, applies a split in one branch, omits it in the other, and
compares their subsequent development-view losses.

The investigation has three hypotheses:
\begin{enumerate}[label=\textbf{H\arabic*:},leftmargin=2.1em]
  \item \textbf{Diagnostic information.} A bounded pose probe changes
  candidate rankings reproducibly beyond zero-update and branch-replay
  variation.
  \item \textbf{Split usefulness.} A response-based feature predicts
  development split utility better than the original gradient and corrected
  gradient magnitude alone, after accounting for candidate attributes and
  visibility.
  \item \textbf{Allocation value.} An allocation policy selected from that
  evidence improves quality versus representation size beyond ordinary pose
  refinement and strong allocation controls, with the complete probe cost
  included.
\end{enumerate}
H1 alone would support a sensitivity diagnostic, not an allocation method.
H2 is required for the proposed predictive mechanism. H3 is required for a
practical reconstruction claim.

The present draft makes no unconditional contribution claim. If supported by
later experiments, the paper may contribute (i) a reproducible measurement of
growth-signal response to a constrained, state-isolated pose intervention,
(ii) direct branch-based evidence about whether that response predicts split
utility, and (iii) an allocation rule whose quality/count trade-off survives
calibration and growth controls. Each item remains \pending{}.

\section{Related Work}

\subsection{Densification and representation allocation}

The original 3DGS method uses an adaptive density-control policy to manage
primitive growth and pruning \cite{kerbl2023}. Revising Densification replaces
parts of this logic with a pixel-error-driven criterion and adds explicit
population-control mechanisms \cite{rota2024}. FastGS uses sampled training
views to score high-error pixels within projected Gaussian footprints, guiding
multi-view-consistent densification and pruning \cite{ren2025fastgs}. These
methods establish that image-error attribution and multi-view evidence are
already central to density control.

Several recent methods further narrow any broad novelty claim about
importance-based or predictive allocation. SalientGS aggregates multi-view
residual evidence into underfit and redundancy signals and uses them to guide
MCMC-based Gaussian allocation in a unified SfM-to-3DGS pipeline
\cite{xiong2026salientgs}. BlitzGS version 2 includes a learned
spawn-survival gate in addition to scheduled population simplification and
workload controls \cite{wang2026blitzgs}. LeGS learns density-control actions
from sampled-view gradient and sensitivity information
\cite{ning2026legs}. Thus, the present question is not whether prediction,
multi-view scoring, or learned allocation can be used in 3DGS. The proposed
distinction is whether the response of an eligible split candidate to a
temporary, bounded extrinsic-pose fit adds information about its subsequent
split utility beyond simpler signals and permanent pose refinement.

\subsection{Pose refinement and nuisance-aware density control}

Energy-GS jointly optimizes scene parameters and camera poses and redesigns
pose-gradient flow with an image-energy-guided alignment procedure
\cite{gao2026energygs}. It is a particularly relevant challenge to a
counterfactual allocation claim: any apparent benefit must be separated from
the benefit of spending more computation on ordinary calibration or changing
the timing of densification.

Other methods address pose or cross-view effects in different observation
settings. BSGS refines camera poses and controls densification for
camera-motion-deblurred images \cite{zhao2025bsgs}. HarmoGS addresses
cross-view gradient conflict in in-the-wild reconstruction and includes
conflict-aware density control \cite{kang2026harmogs}. These works make
pose-related and view-related explanations important controls and failure
cases. Their data assumptions differ from the static, sharp-image setting
specified here.

\subsection{Proposed distinction and its limits}

The targeted review recorded in the Stage 0 artifacts found substantial
overlap in error-driven, multi-view, conflict-aware, learned, and
pose-refinement-based allocation. In the inspected versions, the closest
unresolved question is whether the *change* in a candidate's growth signal
under a small, discarded extrinsic-pose fit predicts its later split benefit,
after accounting for original gradient, corrected magnitude, visibility, and
equivalent permanent pose refinement. The proposed evidence is the conjunction
of isolated intervention measurements, matched split/no-split utility,
and strong calibration controls. This targeted review does not establish
exhaustive novelty. Direct prior art, or a result fully explained by corrected
magnitude or retained pose corrections, would require revising the framing.

Version and reproduction details matter. The local Stage 0 audit checked
version-pinned papers and author-linked repositories on 10 October 2026.
FastGS, SalientGS, BlitzGS, and LeGS have linked source repositories in the
recorded audit. The inspected public default branches for Energy-GS and BSGS
contained README/license files but no implementation suitable for reproduction;
the audit did not identify a runnable HarmoGS repository. These are dated
availability observations, not permanent claims. No external method has been
installed or reproduced for this draft.

\section{Problem Formulation}

Let \(G_t\) denote the Gaussian scene at training step \(t\), \(T_v\) the
camera-to-world transform for training view \(v\), \(K_v\) its fixed
intrinsics, and \(I_v\) its RGB image. Let
\(\mathcal{L}_{\mathrm{photo}}(G,T,K;I)\) be a fixed photometric objective
computed by the shared renderer. The initial study uses
\[
\mathcal{L}_{\mathrm{photo}}
 = 0.8\,\mathcal{L}_{1}
 + 0.2\,(1-\mathrm{SSIM}),
\]
with exact image reduction and SSIM implementation to be recorded with the
training configuration.

For a Gaussian \(i\) eligible for splitting under the common density-control
policy, let \(\boldsymbol{\mu}^{2D}_{iv}\) denote its projected mean in view
\(v\). We use \(g_0(i,v,t)\) and \(g_+(i,v,t)\) for its normalized screen-space
growth-gradient magnitude before and after the temporary pose fit. Under the
repository's DefaultStrategy convention, the two projected-mean gradient
components are scaled by image size and camera count before taking their
norm:
\[
\begin{aligned}
q_x &= \frac{W}{2}\,n_{\mathrm{cam}}\,
       \frac{\partial\mathcal{L}_{\mathrm{photo}}}
            {\partial \mu^{2D}_{iv,x}},\\
q_y &= \frac{H}{2}\,n_{\mathrm{cam}}\,
       \frac{\partial\mathcal{L}_{\mathrm{photo}}}
            {\partial \mu^{2D}_{iv,y}},\\
g(i,v,t) &= \sqrt{q_x^2+q_y^2}.
\end{aligned}
\]
The before/after measurement must use the same image, loss, resolution,
rasterizer flags, Gaussian values, and normalization. A single-view probe
measurement is reported as such; it is not silently equated with the
strategy's accumulated multi-view average.

\section{Proposed Diagnostic}

\subsection{Bounded, state-isolated pose fit}

The intervention operates on an independent copy of a camera. It does not
update Gaussian values, the training camera embedding, optimizer moments,
density-control accumulators, or the main view-sampling sequence. The initial
pose convention is right composition:
\[
T_v^{+}=T_v\Delta_v,\qquad \Delta_v\in\mathrm{SE}(3).
\]
The pose fit minimizes
\[
\begin{split}
\mathcal{L}_{\mathrm{probe}}(\Delta_v)
={}&\mathcal{L}_{\mathrm{photo}}
  (G_t,T_v\Delta_v,K_v;I_v)\\
 &+\lambda_{\mathrm{pose}}\left[
   \left(\frac{\theta(\Delta_v)}{\theta_{\max}}\right)^2
   +\left(
   \frac{\|c(T_v\Delta_v)-c(T_v)\|_2}
        {s_{\mathrm{scene}}\rho_{\max}}
   \right)^2
  \right],
\end{split}
\]
where \(\theta\) is the relative rotation angle, \(c(\cdot)\) is camera center,
and \(s_{\mathrm{scene}}\) is the recorded scene scale. The initial development
setting is five Adam pose-only updates, \(\theta_{\max}=1^\circ\),
\(\rho_{\max}=0.005\), and \(\lambda_{\mathrm{pose}}=0.001\). These are
prospective settings, not empirically selected optima. The bounds apply to
total displacement from the starting camera, not to each update. Any
parameterization or clipping procedure must be logged, and realized camera
motion must be measured in rotation angle and scene-normalized center
displacement.

After fitting the copied pose, the diagnostic recomputes the screen-space
gradient using the photometric term alone. This keeps the before/after
quantities in the same units and prevents the pose prior from directly
changing the gradient comparison. Probe state is discarded. A zero-update
probe and repeated fixed-sample measurements establish the numerical and
replay floor before candidate-ranking changes are interpreted.

\subsection{Candidate validity and score selection}

For every candidate observation, the planned log records candidate identity
and lineage, checkpoint, view, \(g_0\), \(g_+\), photometric loss change,
camera displacement, projected radius, opacity, observation count, and
visibility before and after the intervention. Finite, in-bounds measurements
with stable contribution are analyzed as reliable observations. Disappearing
or newly appearing candidates, large footprint changes, and other visibility
changes are flagged rather than interpreted as calibration-resolved error.
An invalid observation is not negative evidence against splitting; the initial
fallback is the original growth signal, with later reconsideration.

The pilot compares the following rankings:
\begin{center}
\small
\begin{tabularx}{0.96\linewidth}{@{}lY@{}}
\toprule
\textbf{Feature} & \textbf{Question tested}\\
\midrule
Original magnitude \(g_0\) & Does the ordinary growth signal already rank useful splits?\\
Corrected magnitude \(g_+\) & Is one post-correction measurement sufficient?\\
Persistence ratio \(g_+/(g_0+\epsilon)\) & Does relative response add information, despite instability for weak \(g_0\)?\\
Magnitude-aware response \(g_0\min(g_+/(g_0+\epsilon),1)\) & Does combining original magnitude and response help beyond \(g_+\)?\\
Declared model of \(g_0\), \(g_+\), and their difference & Does retaining both measurements add predictive information on held-out candidate batches?\\
Matched random ranking & Does the result exceed suppression with similar gradient, scale, opacity, and visibility statistics?\\
\bottomrule
\end{tabularx}
\end{center}
No final score is frozen at this stage. In particular, the ratio alone is not
treated as a reliable score, and the magnitude-aware example is not claimed as
novel. Stage 3 selects the simplest feature that predicts split utility on
development data; any fitted feature rule is evaluated on candidate batches
or checkpoints not used to fit it.

\section{Experimental Protocol}

\subsection{Scenes and view roles}

Stage 0 prepared two development fixtures and explicit view manifests. The
analytic synthetic scene has known camera truth. The real scene is the
official Mip-NeRF 360 \emph{room} capture, prepared at its low-resolution
images and point-cloud setting \cite{barron2022mipnerf360}.

\begin{table}[htbp]
\centering
\small
\begin{tabularx}{\linewidth}{@{}l l r l Y@{}}
\toprule
\textbf{Scene} & \textbf{Resolution} & \textbf{Points} &
\textbf{Views (train/dev/test)} & \textbf{Role}\\
\midrule
analytic\_spheres\_v1 & \(128\times128\) & 1,280 & 48 (36/6/6) &
Known-camera synthetic development fixture\\
room\_small & \(389\times259\) & 112,627 & 311 (233/39/39) &
Official Mip-NeRF 360 room; previously exposed development scene\\
\bottomrule
\end{tabularx}
\caption{Stage 0 data inventory and prospective three-way view counts.}
\label{tab:data}
\end{table}

For both scenes, view identity, parser index, role, dimensions, and hashes are
recorded in manifests. Training images are used for scene/camera optimization
and probe fitting. Development images are used for score selection and
split-utility branches; probe fitting does not use them in the primary
evaluation. Reserved views are excluded from score and parameter selection.
For reproducibility, if \(i\) is the zero-based lexicographic parser index,
reserved test views satisfy \(i\bmod 8=0\), development views satisfy
\(i\bmod 8=4\), and training views satisfy
\(i\bmod 8\in\{1,2,3,5,6,7\}\).
The room's held-out views have been evaluated in earlier project runs, so the
room is a development fixture rather than untouched final generalization
evidence. Later final evaluation requires fresh scenes and views. The
provided COLMAP initialization may also have used the full capture; all
methods in the initial isolation study share that initialization, and a
strict training-only SfM reconstruction would be a separate setting.

The Stage 0 parser audit verified image decoding, finite camera/point values,
camera-transform consistency, and disjoint manifest roles. Re-running the
analytic-fixture generator reproduced 60 prepared files byte-for-byte. These
are data and artifact checks, not training-quality or method results.

\subsection{Pose perturbations and candidate sampling}

The synthetic scene supports clean-camera evaluation against known pose
metadata. The development plan also includes clean, mild, and moderate
training-pose perturbations. Proposed mild disturbances are at most
\(0.5^\circ\) and \(0.0025\) scene scales; proposed moderate disturbances are
at most \(2^\circ\) and \(0.01\) scene scales. Perturbations are sampled with
fixed seeds and applied to training cameras only. Development and reserved
camera poses remain fixed. These are controlled stress settings, not claims
about typical real-world pose errors.

For the split-utility pilot, the initial plan samples 24 candidates at each of
three checkpoints (25\%, 50\%, and 75\% of the active split window), subject
to measured branch cost. Sampling spans original gradient, scale, opacity,
and observation count; pose-sensitive and persistent candidates are matched
approximately on these attributes. Candidate identities are fixed before
branch outcomes are observed. The initial end-to-end horizons are 3,000 steps
for the analytic fixture and 7,000 for room, with seeds 0, 1, and 2 for
decisive pilot comparisons. Feasibility changes must be recorded before
examining favorable outcomes.

\subsection{Matched split utility}

For each candidate, two branches are restored from the same complete training
snapshot. One branch applies a fixed split to the selected parent; the other
omits it. Other structural changes are held fixed or disabled in both branches
for this short diagnostic. Both branches continue for the same sampled views,
steps, camera policy, optimizer schedule, and state-expansion rule. A
development view set excluded from probe fitting measures utility:
\[
U_i^{(h)}
=
\mathcal{L}_{\mathrm{dev}}(G_{\mathrm{no\ split}}^{(h)})
-
\mathcal{L}_{\mathrm{dev}}(G_{\mathrm{split}\ i}^{(h)}),
\]
where \(h\) is the continuation horizon. Positive \(U_i^{(h)}\) means that
the split branch has lower development loss after the specified continuation.
The initial short horizon is 200 steps, with a prespecified subset continued
for 1,000 steps to assess horizon sensitivity. Grouped splits, if needed for
cost, receive group-level utility labels rather than labels attributed to each
member.

The pilot reports rank correlation, top-quartile utility, complete utility
distributions, invalid-observation rates, and results stratified by relevant
candidate attributes. For 24 candidates, the primary top-quartile size is
\(k=6\). The planned normalized top-\(k\) comparison is
\[
D =
\frac{
  \operatorname{mean}(U_{\mathrm{top}\text{-}k,\mathrm{probe}})
  -\operatorname{mean}(U_{\mathrm{top}\text{-}k,\mathrm{comparator}})
}{
  \operatorname{mean}(\mathcal{L}_{\mathrm{dev,no\ split}})
}.
\]
The prospective pilot signal criterion is \(D\geq0.005\), with the raw
utility difference also exceeding five times the 95th-percentile absolute
no-intervention branch-loss difference measured on matched replays. A
directional advantage over both \(g_0\) and \(g_+\) is required at two
checkpoints, with a second-fixture check and no systematic long-horizon
reversal. This is a development decision rule, not a significance test or a
claim of statistical power.

\subsection{State control and replay}

Reliable branch comparisons require snapshots beyond a scene-only checkpoint.
The planned snapshot contains Gaussian tensors, actual camera state, optimizer
and scheduler states, density-control accumulators, training step, candidate
identities, Python/NumPy/PyTorch CPU and CUDA random-number states, and the
subsequent view-sampling sequence. Appearance state is included if enabled.
Two no-intervention continuations from one snapshot establish replay
agreement within a documented numerical tolerance. The current trainer's
scene checkpoint does not provide this full-state replay and its checkpoint
path is evaluation-only; Stage 1 must add a separate snapshot/restore path
before split effects are measured.

\subsection{Comparisons and resource accounting}

The primary practical comparison is ordinary density control with the same
real camera-pose optimizer. The method must also face controls that isolate
extra calibration work, use of the corrected pose, simpler scores, and mere
growth reduction.

\begin{table}[htbp]
\centering
\small
\begin{tabularx}{\linewidth}{@{}lY@{}}
\toprule
\textbf{Comparison} & \textbf{Alternative explanation tested}\\
\midrule
Ordinary density control & Reference training policy\\
Ordinary density control plus pose refinement & Whether pose optimization alone explains a gain\\
Equal-computation pose-only updates & Whether additional calibration compute explains a gain\\
Retain the probe correction and densify normally & Whether accepting the correction is sufficient\\
Corrected-gradient magnitude \(g_+\) & Whether the before/after response adds value beyond \(g_+\)\\
Multi-view scoring without the pose probe & Whether view consistency explains the allocation\\
Delayed densification & Whether waiting for a more mature scene is sufficient\\
Higher threshold or matched random suppression & Whether simply reducing growth explains the result\\
FastGS and other selected external methods & Whether the result survives the closest available published allocation baselines\\
\bottomrule
\end{tabularx}
\caption{Essential internal controls and external-baseline priorities.}
\label{tab:controls}
\end{table}

Quality versus Gaussian count is the primary end-to-end outcome. At a declared
development operating point, the prospective comparable-quality margins are
\[
\Delta\mathrm{PSNR}\geq-0.20\ \mathrm{dB},\qquad
\Delta\mathrm{SSIM}\geq-0.005,\qquad
\Delta\mathrm{LPIPS}\leq+0.010.
\]
The initial compactness aspiration is no more than 80\% of the reference
final Gaussian count at comparable quality. These values were fixed before
CCD-GS experiments; they are aspirations, not predicted benefits or
publication thresholds. The method must be compared to pose-enabled ordinary
training and to retained-correction/equal-computation controls, not only to
vanilla fixed-pose training.

For every run, report per-scene PSNR, SSIM, and LPIPS; final and peak Gaussian
counts; model size; peak VRAM; total wall-clock time; probe rendering,
backpropagation, scoring, and allocation overhead; and pose error where
reference poses exist. Fixed-budget and fixed-total-time comparisons are
separate: a fixed number of steps does not guarantee equal computation.
Rendering speed, if studied, requires its own synchronized fixed-hardware
benchmark. Primitive count alone is not rendering-speed evidence.

The later frozen campaign targets at least six fresh scenes across two
dataset families, including known-camera synthetic scenes and ordinary real
reconstructions. FastGS is the first external allocation baseline to
investigate. Protocol-compatible reproduction, hardware/software, renderer
changes, and unavailable or ported methods must be described explicitly.

\section{Current Evidence and Results Status}

Stage 0 completed the question, scope, targeted related-work review, data
preparation, view manifests, and prospective decision criteria on 10 October
2026. It did not execute baseline training, pose-probe optimization,
split/no-split branches, integrated allocation, external method reproduction,
or final test scoring.

\begin{table}[htbp]
\centering
\small
\begin{tabularx}{\linewidth}{@{}l l Y@{}}
\toprule
\textbf{Claim under test} & \textbf{Status} & \textbf{Required evidence}\\
\midrule
Probe changes rankings reproducibly (H1) & Not measured &
Repeated fixed-state probes, zero-update control, replay and visibility analysis\\
Response predicts split utility (H2) & Not measured &
Matched split/no-split branches; comparison with \(g_0\), \(g_+\), and matched random rankings\\
Allocation improves quality/count (H3) & Not measured &
Multi-budget curves and strong pose/growth controls across fresh scenes\\
Training-time improvement & Not measured &
Equal-total-time accounting that includes all probe costs\\
Rendering-time improvement & Not measured &
Independent synchronized rendering benchmark\\
Camera-pose improvement & Not measured &
Known-pose evaluation under a fixed gauge policy\\
\bottomrule
\end{tabularx}
\caption{Evidence ledger at the time of this draft.}
\label{tab:evidence}
\end{table}

The completed data checks establish that the selected development fixtures
and view-role manifests can be inspected and loaded through the parser. They
do not establish that the baseline trains under the prospective three-way
split, that matched snapshots replay reliably, or that any probe signal
predicts useful growth. No CCD-GS numerical result is reported.

\section{Discussion}

The proposed mechanism is deliberately falsifiable. If pose-probe rankings
are not reproducible beyond replay and numerical variation, H1 fails. If
corrected-gradient magnitude predicts utility equally well, a more complicated
before/after score is not justified. If retaining probe corrections or
spending the same time on ordinary pose updates matches the result, the
evidence would support calibration or scheduling effects rather than a
counterfactual allocation advantage. If the score predicts split utility but
does not improve end-to-end quality/count trade-offs, the paper should remain
a diagnostic study. A reliable negative result is also reportable.

Several failure modes need explicit analysis. Pose-scene ambiguity can make
camera movement hide geometric error. Visibility changes can alter a gradient
without resolving the scene. Persistent residuals may indicate that existing
Gaussians are under-optimized rather than that the representation needs more
capacity. A temporary correction that is discarded may not be reached by the
ordinary optimizer. Finally, permanently suppressing candidates can starve
regions while the scene is immature. These risks motivate copied probe state,
visibility logging, invalid-observation fallback, later reconsideration, and
the retained-correction control.

This draft also has limits before any method results exist. The literature
review is targeted rather than exhaustive. The synthetic fixture is small and
analytic. The room scene is previously exposed development data and cannot
serve as untouched generalization evidence. Its supplied camera/point-cloud
initialization may reflect the full capture. Split utility depends on the
chosen continuation horizon and candidate interactions. The planned early
pilot can reject an implausible mechanism, but it cannot replace a frozen
multi-scene campaign with repeated runs and fresh final views.

\section{Conclusion}

We have specified a study of whether a bounded, temporary camera-extrinsic
intervention helps predict the utility of Gaussian splits. Its central test is
not whether gradients change after a pose fit, but whether that change
predicts held-out development split utility and improves quality/count
allocation beyond ordinary pose refinement and equal-computation controls.
Stage 0 prepared development data and recorded prospective decision rules;
it produced no method-performance evidence. The question therefore remains
open, and the final paper claims must follow the results of the planned
experiments.

\section*{Acknowledgments}

\textit{To be completed.}

\begin{thebibliography}{99}
\small

\bibitem{kerbl2023}
B. Kerbl, G. Kopanas, T. Leimk\"uhler, and G. Drettakis.
\newblock 3D Gaussian Splatting for Real-Time Radiance Field Rendering.
\newblock \emph{ACM Transactions on Graphics}, 42(4), 2023.
\newblock \url{https://arxiv.org/abs/2308.04079}.

\bibitem{rota2024}
S. Rota Bul\`o, L. Porzi, and P. Kontschieder.
\newblock Revising Densification in Gaussian Splatting.
\newblock In \emph{European Conference on Computer Vision (ECCV)}, 2024.
\newblock \url{https://arxiv.org/abs/2404.06109}.

\bibitem{ren2025fastgs}
S. Ren, T. Wen, Y. Fang, and B. Lu.
\newblock FastGS: Training 3D Gaussian Splatting in 100 Seconds.
\newblock In \emph{Proceedings of the IEEE/CVF Conference on Computer Vision
and Pattern Recognition (CVPR)}, 2026. arXiv:2511.04283, version 3.
\newblock \url{https://arxiv.org/abs/2511.04283}.

\bibitem{xiong2026salientgs}
T. Xiong, R. Li, S. Ge, and J. Yang.
\newblock SalientGS: Unified SfM-to-3DGS with Importance-Guided MCMC Gaussian
Allocation.
\newblock In \emph{Proceedings of the ACM International Conference on
Multimedia}, 2026. arXiv:2607.11285, version 2.
\newblock \url{https://arxiv.org/abs/2607.11285}.

\bibitem{gao2026energygs}
Y. Gao, L. Su, R. Huang, T. Jiang, J. Tang, Y. Yue, and Y. Yang.
\newblock Energy-GS: Image Energy-guided Pose Alignment Gaussian Splatting
with redesigned pose gradient flow.
\newblock In \emph{Proceedings of the IEEE/CVF Conference on Computer Vision
and Pattern Recognition (CVPR)}, pages 7310--7319, 2026.
\newblock \url{https://openaccess.thecvf.com/content/CVPR2026/html/Gao_Energy-GS_Image_Energy-guided_Pose_Alignment_Gaussian_Splatting_with_redesigned_pose_CVPR_2026_paper.html}.

\bibitem{zhao2025bsgs}
A. Zhao, P. Yu, Z. Zhu, and M. Wei.
\newblock BSGS: Bi-stage 3D Gaussian Splatting for Camera Motion Deblurring.
\newblock In \emph{Proceedings of the ACM International Conference on
Multimedia}, 2025. arXiv:2510.12493, version 2.
\newblock \url{https://arxiv.org/abs/2510.12493}.

\bibitem{kang2026harmogs}
Y. Kang, T. Zhu, J.-F. Hu, J. Lai, and W.-S. Zheng.
\newblock HarmoGS: Robust 3D Gaussian Splatting in the Wild via
Conflict-Aware Gradient Harmonization.
\newblock arXiv:2605.13073, version 2, 2026.
\newblock \url{https://arxiv.org/abs/2605.13073}.

\bibitem{wang2026blitzgs}
Z. Wang, H. Au, Y. Li, M. Su, H. Jin, Y. Chen, M. Gai, F. Zhu, and G. Wang.
\newblock BlitzGS: City-Scale Gaussian Splatting at Lightning Speed.
\newblock arXiv:2605.13794, version 2, 2026.
\newblock \url{https://arxiv.org/abs/2605.13794}.

\bibitem{ning2026legs}
Z. Ning, X. Li, J. Yu, G. Lu, Y. Wang, and W. Pei.
\newblock Beyond Heuristics: Learnable Density Control for 3D Gaussian
Splatting.
\newblock arXiv:2605.00408, version 2, 2026.
\newblock \url{https://arxiv.org/abs/2605.00408}.

\bibitem{barron2022mipnerf360}
J. T. Barron, B. Mildenhall, D. Verbin, P. P. Srinivasan, and P. Hedman.
\newblock Mip-NeRF 360: Unbounded Anti-Aliased Neural Radiance Fields.
\newblock In \emph{Proceedings of the IEEE/CVF Conference on Computer Vision
and Pattern Recognition (CVPR)}, 2022.
\newblock \url{https://jonbarron.info/mipnerf360/}.

\end{thebibliography}
\end{document}
